\part{Appearance, labels and evaluation}

\section{Appearance, reusable styles and palettes}
Ten themes are available: blue, teal, solid, gray, lbyellow, viblue, cyblu,
violet, orange and forest. The three new themes pair violet, warm-orange or
forest-green gradients with matching pale tracks and dark value labels.
Nine predefined solid colors are lightgreen, lightyellow, lightblue, lightred,
rose, skyblue, gold, lavender and peach. Use \key{color} for solid fill and
\key{left color}/\key{right color} for a gradient. Standard xcolor expressions work.

\begin{demo}{Violet, warm orange and forest green themes}
\gradbar[theme=violet]{72}

\gradbar[theme=orange]{72}

\gradbar[theme=forest]{72}
\end{demo}

\begin{demo}{A reusable layout style}
\gradbarsstyle{paperrow}{preset=paper,width=45mm,max=100,precision=0}
\gradbar[style=paperrow]{82}

\graddumbbell[style=paperrow]{60}{82}

\gradbar[preset=report,target=90,thresholds={60,80}]{75}

\gradbar[preset=presentation]{92}

\gradbar[preset=outline]{70}
\end{demo}
The paper preset uses thin solid bars, a faint track and outside values. Report
uses thicker rounded bars; supply actual target and threshold values yourself.
Presentation uses larger type and automatic inside placement. Outline uses hollow
shapes with fine strokes. Presets do not invent data or scale limits.

Group palettes are categorical (six distinct colors), sequential (five blue levels),
diverging (orange through neutral to teal), and mono (gray with textures).
These palettes select by segment position, not by value; short lists cycle.
Named categories instead preserve identity across reordered or subsetted data:
\begin{demo}{Stable category identity and a shared legend}
\gradbarscategory{Compute}{gradbarsBlue}{diagonal}
\gradbarscategory{Storage}{gradbarsOrange}{dots}
\gradstack[stack names={Compute,Storage}]{60,30}

\gradstack[stack names={Storage,Compute}]{30,60}

\gradbarslegend[width=50mm]{Compute,Storage}
\end{demo}
Registered names override positional colors and patterns; unregistered names fall
back to the current palette. Definitions are scoped. In print/mono mode, category
colors use gray while registered patterns remain. Legends use the same mapping.
Patterns are none, diagonal, reverse, dots, crosshatch, horizontal and vertical.
Use names or textures alongside color when category distinctions matter.
